\documentclass[11pt]{article}
\usepackage{amsmath,amssymb,amsthm}

\newtheorem{theorem}{Theorem}
\newtheorem{lemma}[theorem]{Lemma}
\newtheorem{corollary}[theorem]{Corollary}

\title{A terminal-tower reduction for high Collatz decoder depth}
\author{}
\date{27 August 2026}

\begin{document}
\maketitle

Let
\[
  2^{a_i}x_{i+1}=3x_i+1,
  \qquad
  a_i=v_2(3x_i+1),
  \tag{1}
\]
where the $x_i$ are positive odd integers.  Put
\[
  A_k:=\sum_{i=0}^{k-1}a_i,
  \qquad
  b_i:=v_3(x_i+1),
  \qquad
  x_i+1=3^{b_i}u_i,
  \qquad
  3\nmid u_i.
  \tag{2}
\]
We call the exponent block $(a_0,\ldots,a_{r-1})$ a strict barrier block if
\[
  2^{A_k}<3^k
  \qquad(1\leq k\leq r).
  \tag{3}
\]
Set $\alpha:=\log_2 3$.  For $a\geq1$, define
\[
  q(a):=v_3(2^a-2),
  \qquad q(1):=\infty.
  \tag{4}
\]

\begin{lemma}[one-step source dichotomy]
Fix one step, and write $a:=a_i$, $b:=b_i$, and $b':=b_{i+1}$.  If
$q(a)\neq b+1$, then
\[
  b'=\min\{q(a),b+1\}.
  \tag{5}
\]
If $q(a)=b+1<\infty$, put
\[
  \rho_a:=\frac{2^a-2}{3^{q(a)}}.
  \tag{6}
\]
Then
\[
  b'=b+1+v_3(u_i+\rho_a).
  \tag{7}
\]
Moreover,
\[
  q(a)=
  \begin{cases}
    0,&a\text{ is even},\\
    1+v_3((a-1)/2),&a>1\text{ is odd}.
  \end{cases}
  \tag{8}
\]
\end{lemma}

\begin{proof}
Adding $2^a-1$ to (1) gives
\[
  2^a(x_{i+1}+1)=3(x_i+1)+2^a-2.
  \tag{9}
\]
The two terms on the right have valuations $b+1$ and $q(a)$.  Distinct
valuations give (5), while equal finite valuations give (7) after division
by $3^{b+1}$.

If $a$ is even, then $2^a-2\not\equiv0\pmod3$.  If $a=2n+1>1$, then
\[
  2^a-2=2(4^n-1).
\]
Writing $n=3^s m$ with $3\nmid m$, the congruence
$(1+3)^m-1\equiv3m\pmod9$ gives $v_3(4^m-1)=1$.  The factorization
$X^3-1=(X-1)(X^2+X+1)$, with
$X^2+X+1\equiv3\pmod9$ for $X\equiv1\pmod3$, raises the valuation by one
whenever the exponent is multiplied by $3$.  This proves (8).
\end{proof}

\begin{theorem}[terminal tower or terminal collision]
Let $(a_0,\ldots,a_{r-1})$ be a strict barrier block, and define
\[
  K(r):=\left\lceil\log_3\left(\frac{\alpha r}{2}\right)\right\rceil.
  \tag{10}
\]
Let $\ell$ be the length of its terminal run of exponent-one letters.  If
$\ell=r$, then
\[
  b_r=b_0+r.
  \tag{11}
\]
If $\ell<r$, put $s:=r-\ell-1$, so $a_s>1$, and define
\[
  g_s:=
  \begin{cases}
    v_3(u_s+\rho_{a_s}),&q(a_s)=b_s+1,\\
    0,&q(a_s)\neq b_s+1.
  \end{cases}
  \tag{12}
\]
Then
\[
\boxed{
  b_r\leq K(r)+\ell+g_s.}
\tag{13}
\]
Thus every decoder depth above the logarithmic scale $K(r)$ comes from the
terminal exponent-one tower, from the final non-one collision, or from both.
\end{theorem}

\begin{proof}
Condition (3) gives
\[
  a_i\leq A_{i+1}<\alpha(i+1)\leq\alpha r.
  \tag{14}
\]
For every odd $a_i>1$, equation (8) gives
\[
  3^{q(a_i)-1}
  \leq\frac{a_i-1}{2}
  <\frac{\alpha r}{2}
  \leq3^{K(r)}.
\]
Hence $q(a_i)\leq K(r)$.  The same bound holds for even $a_i$, since then
$q(a_i)=0$.

Suppose $\ell<r$.  If the step at $s$ is not a collision, (5) gives
\[
  b_{s+1}\leq q(a_s)\leq K(r).
  \tag{15}
\]
If it is a collision, (7) and $q(a_s)=b_s+1$ give
\[
  b_{s+1}=q(a_s)+g_s\leq K(r)+g_s.
  \tag{16}
\]
Every one of the remaining $\ell$ exponent-one steps raises the decoder
depth by exactly one.  Equations (15) and (16) therefore prove (13).  The
same one-step law proves (11).
\end{proof}

\begin{corollary}[deficit-three rate reduction]
Under the hypotheses of Theorem~2, suppose that the terminal decoder has
deficit at most three:
\[
  b_r\geq r-3.
  \tag{17}
\]
If the block is not all ones, then
\[
\boxed{
  \ell+g_s
  \geq r-3-K(r).}
\tag{18}
\]
In particular, for every $0\leq\eta<1$,
\[
  \ell\leq\eta r
  \quad\Longrightarrow\quad
  g_s\geq(1-\eta)r-3-K(r).
  \tag{19}
\]
Whenever the right side of (19) is positive, the last non-one step obeys the
explicit collision conditions
\[
  q(a_s)=b_s+1,
  \qquad
  u_s\equiv-\rho_{a_s}\pmod {3^{g_s}}.
  \tag{20}
\]
Consequently, along any family of strict barrier blocks satisfying (17), a
sublinear terminal exponent-one tower forces a collision of depth
$r-o(r)$ at the last non-one step.
\end{corollary}

\begin{proof}
Equation (18) is (13) rearranged under (17), and (19) follows immediately.
If its right side is positive, then $g_s>0$.  Definition (12) gives both
claims in (20).  Since $K(r)=O(\log r)$, the final assertion follows from
(19).
\end{proof}

\end{document}
